\section{A Taxonomy of Adversarial Pragmatics}

\term{Adversarial pragmatics} names safety-relevant cases in which the operative failure mode depends on how language is embedded, attributed, scoped, or taken up. The taxonomy is intentionally instrumental: it exists to improve eval construction and failure analysis, not to settle a general theory of meaning.

The initial benchmark distinguishes eight families.

\subsection{Embedded commands}

The model has to distinguish the user's instruction from text embedded in a webpage, email, document, log, table, code comment, citation, or tool output.

\subsection{Mention/use and quotation}

The safe action depends on whether a policy-relevant string is being mentioned, classified, translated, reported, or enacted.

\subsection{Authority and instruction hierarchy}

The model has to resolve conflicts among system, developer, user, tool, document, and external-source instructions.

\subsection{Scope, negation, and modality}

Words such as \mention{not}, \mention{unless}, \mention{only if}, \mention{may}, \mention{must}, \mention{should}, and \mention{except} change what counts as compliance.

\subsection{Deixis and reference hijacking}

Indexicals and anaphora such as \mention{this}, \mention{that}, \mention{the above}, \mention{the previous instruction}, \mention{the user}, and \mention{the admin} can be redirected.

\subsection{Indirect speech acts and pragmatic pressure}

Requests may be framed as hypotheticals, jokes, roleplay, bureaucratic boilerplate, urgency, politeness, or \enquote{just checking}.

\subsection{Policy-boundary ambiguity}

The case isn't simply safe or unsafe; the policy construct itself may be underspecified, unstable across paraphrase, or sensitive to rater role.

\subsection{Agent transcript interpretation}

The evaluator has to decide whether a transcript shows refusal, inability, tool failure, scaffold failure, instruction conflict, goal drift, omitted information, or strategic misreporting.
